\documentclass[11pt,a4paper]{article}
\usepackage[utf8]{inputenc}
\usepackage[ngerman]{babel}
\usepackage{graphicx}
\usepackage{cite}
\usepackage{geometry}
\usepackage{booktabs}
\usepackage{array}
\usepackage{url}
\usepackage{amsmath}
\usepackage{microtype}
\usepackage{parskip}
\usepackage{tikz}
\usetikzlibrary{arrows.meta, positioning}
\geometry{margin=1in}
\emergencystretch=1em
\setlength{\abovecaptionskip}{5pt}
\setlength{\belowcaptionskip}{2pt}
\setlength{\textfloatsep}{10pt}

\title{Entwicklung einer mehrstufigen Pipeline zur autonomen
Erstellung wissenschaftlicher Paper:\\
Architekturentscheidungen, Beobachtungen und Evaluation}
\author{Jakob Gl\"uck \\[4pt]
\small Technische Universit\"at Chemnitz \\
\small Hauptseminar Medieninformatik 2026}
\date{Abgabe: 17.07.2026}

\begin{document}

\maketitle

\begin{abstract}
Dieser Bericht dokumentiert die iterative Entwicklung einer mehrstufigen KI-Pipeline
zur autonomen Generierung wissenschaftlicher Arbeiten. Ausgangspunkt war ein
unstrukturierter Einzelprompt ohne Werkzeugzugang, die Baseline. Die Ergebnisse
wiesen erhebliche methodische Mängel auf. Über drei Iterationen (v1 bis v3) wurden
gezielte Verbesserungen eingeführt: Werkzeuganbindung, interne Kontrollmechanismen,
Pre-Registrierung und schließlich ein adversarialer Critic-Prozess, der jede Version
auf die Schwachstellen der vorherigen prüfte. Die Evaluation belegt eine deutliche
Qualitätssteigerung: von 9 von 25 Punkten (Baseline) auf 24 von 25 Punkte in
Version~3. Ein zentraler Befund replizierte sich über alle Iterationen stabil: der
negative Zusammenhang zwischen wahrgenommener KI-Jobbedrohung und Jobzufriedenheit
($d$ liegt zwischen $-0{,}30$ und $-0{,}36$).
\end{abstract}

\section{Einleitung}

Sprachmodelle verfassen komplexe Texte, generieren Code und führen statistische Analysen
durch. Ob sie damit auch den gesamten Forschungsprozess autonom durchlaufen können, von der
Datenexploration bis zum fertigen Paper, ist bislang ungeklärt.

Das vorliegende Projekt untersucht diese Fragestellung anhand einer inkrementell
entwickelten Agenten-Pipeline. Als empirische Datenbasis diente der Stack Overflow
Developer Survey 2024 \cite{stackoverflow2024} mit rund 65.000 Teilnehmenden weltweit.
Die Pipeline agiert vollständig autonom: Nach Übergabe des Datensatzes formuliert das
System selbstständig Forschungsfragen, führt die statistische Auswertung durch und
verfasst das wissenschaftliche Dokument. Vergleichbare Frameworks existieren bereits,
etwa AI~Scientist \cite{lu2024aiscientist} oder Agent~Laboratory
\cite{schmidgall2025agentlab}. Der Unterschied liegt nicht in maximaler Automatisierung, sondern in strukturierter
Qualitätssicherung.

Die Baseline-Untersuchung zeigte fundamentale Limitationen unstrukturierter Prompts.
Ohne Echtzeitzugriff auf externe Werkzeuge neigen Modelle zu Halluzinationen bei
Literaturangaben \cite{rao2026detecting}. LLMs übersehen bei der Selbstüberprüfung 52 bis 82\,\% ihrer eigenen methodischen
Fehler \cite{badscientist2025}. Zudem tendieren sie dazu, statistische Signifikanz
mit praktischer Relevanz gleichzusetzen \cite{badscientist2025}.

Um diesen Defiziten zu begegnen, wurde die Pipeline in drei Stufen ausgebaut. Qualität wird dabei verstanden als die Erfüllung verifizierbarer Kriterien:
valide Zitationen, methodische Robustheit, Schutz vor HARKing und Reproduzierbarkeit. Forschungsfrage und Hypothesen lauteten:

\begin{description}
  \item[\textbf{FF:}] Verbessert eine strukturierte Pipeline die wissenschaftliche
    Qualität im Vergleich zu einem unstrukturierten Prompt?
  \item[$H_0$:] Eine strukturierte Pipeline erzeugt gegenüber einem unstrukturierten Prompt keinen messbaren Qualitätsunterschied.
  \item[$H_1$:] Eine strukturierte Pipeline erzeugt gegenüber einem unstrukturierten Prompt messbare Qualitätsverbesserungen.
\end{description}

\section{Grundlagen}

\subsection{Große Sprachmodelle}

Große Sprachmodelle sind neuronale Netze, trainiert auf umfangreichen Textsammlungen
\cite{brown2020gpt3}. Gesteuert werden sie über Prompts; die Ausgabe basiert auf
statistischen Mustern der Trainingsdaten. Für den wissenschaftlichen Einsatz ist relevant, dass moderne Modelle nicht nur
natürliche Sprache verfassen, sondern auch Programmcode sowie strukturierte
Datenformate erzeugen.

\subsection{Autonome Agenten und Pipelines}

Ein Chatbot reagiert auf Eingaben und antwortet. Ein autonomer Agent geht darüber
hinaus: Er erhält ein Ziel und entscheidet eigenständig, welche Schritte dorthin
führen \cite{lu2024aiscientist, schmidgall2025agentlab}. In diesem Projekt erhält
der Agent einen Datensatz und liefert am Ende ein Paper ab, ohne menschliche
Eingriffe. Eine Pipeline gibt diesem Prozess eine Struktur, teilt ihn in feste
Phasen und steuert,
wie Informationen zwischen diesen Phasen fließen.

\subsection{Model Context Protocol (MCP)}

Das Model Context Protocol (MCP) \cite{anthropic2025mcp} gibt Sprachmodellen Zugriff
auf externe Werkzeuge. Über MCP-Server kann ein Modell Datenbanken abfragen,
Webseiten abrufen oder Dateien bearbeiten. Zwei
MCP-Werkzeuge spielen in diesem Projekt eine zentrale Rolle: SQLite-MCP für den
strukturierten Datenbankzugriff und Fetch-MCP für den Live-Abruf von Literaturquellen.

\section{Datensatz und experimentelles Design}

\subsection{Datensatz}

Als Datengrundlage diente der \emph{Stack Overflow Developer Survey 2024}
\cite{stackoverflow2024}, eine jährliche Befragung von Softwareentwicklerinnen
und -entwicklern weltweit ($N = 65{,}437$, 114 Variablen). Die relevanten
Variablen für dieses Projekt sind die wahrgenommene KI-Jobbedrohung
(\texttt{AIThreat}), das Vertrauen in KI-Ausgaben (\texttt{AIAcc}), der
KI-Nutzungsstatus (\texttt{AISelect}) sowie die Jobzufriedenheit
(\texttt{JobSat}, 0 bis 10). Die Variable \texttt{JobSat} weist einen Missing-Anteil
von 55{,}5\,\% auf, da die Frage nur an Berufstätige gerichtet war.

\subsection{Versuchsdesign}

Es wurden vier Bedingungen verglichen, die strukturell aufeinander aufbauen.
Die \textbf{Baseline} entspricht einem unkontrollierten Einzelprompt ohne
Werkzeugzugang. Die \textbf{Versionen v1 bis v3} sind mehrstufige Pipelines,
die in Architektur und eingesetzten Mechanismen schrittweise erweitert wurden.
In allen vier Bedingungen wurden derselbe Datensatz und dasselbe Sprachmodell
(Claude Sonnet 4.6) eingesetzt, ohne menschliche Eingriffe während der Ausführung.

\begin{table}[h]
\centering
\caption{Strukturmerkmale der vier Bedingungen}
\label{tab:design}
\small
\begin{tabular}{lcccc}
\toprule
\textbf{Merkmal} & \textbf{Baseline} & \textbf{v1} & \textbf{v2} & \textbf{v3} \\
\midrule
Stufenarchitektur & keine & 5 Stages & 5 Stages & 7 Stages \\
Datenbankzugriff & Nein & SQLite-MCP & SQLite-MCP & SQLite-MCP \\
Literaturverifikation & Nein & Fetch-MCP & Fetch-MCP und Critic & PRISMA und Critic \\
QC mit STDOUT-Beleg & Nein & Ja & Ja & Ja \\
Pre-Registrierung & Nein & Nein & Nein & Ja \\
Adversarial Critic & Nein & Nein & Nein & Ja \\
\bottomrule
\end{tabular}
\end{table}

\section{Pipeline-Architektur und Designentscheidungen}

\subsection{Baseline: Grenzen des unkontrollierten Prompts}

Unter der Baseline-Bedingung formulierte der Agent aus einem einzigen Prompt eine
Forschungsfrage, analysierte Daten und erstellte ein Paper. Ohne Werkzeugzugriff traten gravierende Defizite auf: Literaturangaben konnten nicht
in Echtzeit verifiziert werden, die Herkunft berechneter Werte ließ sich nicht
nachvollziehen, und die Korrektheit von Seitenzahlen war nicht kontrollierbar. Diese Defizite
bildeten die Grundlage für die Architekturentscheidungen in v1.

\subsection{Pipeline v1: Werkzeugintegration und Stufentrennung}

Die Kernidee von v1 war eine klare Aufteilung in Stufen. Datenerkundung,
Literaturrecherche, Analyse, Schreiben und Qualitätskontrolle liefen sequenziell.
Das verhinderte, dass das Modell Exploration und Ergebnisdarstellung vermischte.

Ergänzt wurde die Architektur durch zwei MCP-Werkzeuge: SQLite-MCP für direkten
Datenbankzugriff und Fetch-MCP zum Echtzeit-Abruf von Literaturquellen. MCP \cite{anthropic2025mcp} erwies sich gegenüber Websuchanfragen als überlegen,
da der Agent jeden Abstract lesen und die Relevanz dokumentieren musste, bevor
eine Quelle zitiert werden durfte.

In den initialen Testläufen wurden drei kritische Schwachstellen identifiziert und
behoben. Erstens übernahm der Agent eine E-Mail-Adresse aus dem Systemkontext in
das Autorenfeld. Textbasierte Anweisungen erwiesen sich als unwirksam; als Lösung
wurde eine deklarative Anonymous-Author-Vorgabe in allen Stages verankert. Zweitens
dokumentierte der Agent erfolgreiche Qualitätskontroll-Durchläufe (QC), ohne das
Prüfskript tatsächlich auszuführen. Als Gegenmaßnahme wurde eine STDOUT-Pflicht implementiert: Als gültiger
PASS wird ausschließlich verifizierter Terminal-Output gewertet. Das reduziert das
Risiko, schließt jedoch nicht aus, dass ein Modell STDOUT fabrizieren kann. Drittens
wählte der Agent Forschungsfragen mit praktisch vernachlässigbaren Effekten. Ein
Effect-Size-Gate ($|d| \geq 0{,}25$) verpflichtete das System zur Suche nach
praktisch bedeutsamen Effekten.

Besonders relevant war die sogenannte "`Likert-Rebellion"': Der Agent
ignorierte eine explizite Dummy-Coding-Vorgabe und behandelte die Variablen
stattdessen ordinal. Die Begründung des Agenten war statistisch: Dummy-Coding ignoriert die Rangordnung
der Likert-Kategorien. Leistungsfähige Modelle wägen Textanweisungen gegen ihr
methodisches Wissen ab, weshalb methodische Vorgaben im Code verankert sein sollten
und nicht ausschließlich im Prompt.

\subsection{Pipeline v2: Interne Verifikationsschleifen}

Während v1 Verifikation vorsah, fehlte eine aktive Rückkopplung. v2 integrierte
daher Critic-Loops in jede Phase. Nach jeder Fetch-MCP-Abfrage wurde jede Quelle ein
zweites Mal abgerufen und inhaltlich überprüft. Für zentrale Zahlen gab es eine
unabhängige Nachrechnung per SQL, die Abweichungen als Fehler markierte. In
Stage~3 wurde vor jedem Satz mit Quellenangabe die URL erneut abgerufen und der
Inhalt geprüft.

Diese Verifikationsschleifen erwiesen sich als wirksam. Zwei hinter Paywalls gesperrte
Quellen wurden durch frei zugängliche ersetzt. Ein Multikollinearitätsfehler im
Regressionsmodell wurde erkannt und als PAP-Abweichung dokumentiert. Eine überspitzte Aussage im Diskussionsteil wurde korrigiert. In allen drei Fällen
deckten die Schleifen reale Fehler auf, die ohne strukturierte Verifikation
unbemerkt geblieben wären.

v2 schützte vor Messfehlern, aber nicht vor systematischen Verzerrungen im
Forschungsdesign selbst. v3 adressierte das mit zwei grundlegenden Erweiterungen.

\subsection{Pipeline v3: Pre-Registrierung und Adversarial Critic}

\begin{figure}[ht]
\centering
\begin{tikzpicture}[
  node distance=0.35cm and 1.8cm,
  box/.style={rectangle, draw=black, rounded corners=3pt,
              minimum width=5.2cm, minimum height=0.65cm,
              text centered, font=\small},
  newbox/.style={box, fill=black!12, draw=black!70},
  arr/.style={-{Stealth[length=5pt]}, thick},
  note/.style={font=\footnotesize\itshape, text=black!55, anchor=west}
]
\node[box]    (s0)  {Stage~0: Datenexploration};
\node[newbox, below=of s0]  (s0b) {Stage~0b: Pre-Analysis Plan};
\node[box,    below=of s0b] (s1)  {Stage~1: Literaturrecherche};
\node[box,    below=of s1]  (s2)  {Stage~2: Statistische Analyse};
\node[newbox, below=of s2]  (s2b) {Stage~2b: Adversarial Critic};
\node[box,    below=of s2b] (s3)  {Stage~3: Paper-Verfassung};
\node[box,    below=of s3]  (s4)  {Stage~4: Qualitätskontrolle};
\draw[arr] (s0)  -- (s0b);
\draw[arr] (s0b) -- (s1);
\draw[arr] (s1)  -- (s2);
\draw[arr] (s2)  -- (s2b);
\draw[arr] (s2b) -- (s3);
\draw[arr] (s3)  -- (s4);
\node[note, right=0.5cm of s0]  {v1, v2, v3};
\node[note, right=0.5cm of s0b] {neu in v3};
\node[note, right=0.5cm of s1]  {v1, v2, v3};
\node[note, right=0.5cm of s2]  {v1, v2, v3};
\node[note, right=0.5cm of s2b] {neu in v3};
\node[note, right=0.5cm of s3]  {v1, v2, v3};
\node[note, right=0.5cm of s4]  {v1, v2, v3};
\end{tikzpicture}
\caption{Stufenarchitektur der Pipeline in der vollständigen Version v3. Grau hinterlegt: in v3 neu eingeführte Stages; v1 und v2 umfassen Stage~0 bis Stage~4 ohne Stage~0b und~2b.}
\label{fig:pipeline}
\end{figure}

\textbf{Pre-Analysis Plan (Stage~0b):} Bevor der Agent Zugriff auf Daten erhielt,
wurde eine vollständige Präregistrierung durchgeführt. Hypothese, statistische
Methode, Mindest-Effektgröße und ein kausaler DAG wurden als JSON-Artefakt gesperrt.
Alle späteren Abweichungen mussten explizit als exploratorisch gekennzeichnet werden.
HARKing wurde damit strukturell unmöglich gemacht, nicht nur per Textanweisung. Das
orientiert sich an Registered Reports \cite{scheel2021excess}.

\textbf{Adversarial Critic (Stage~2b):} Der Agent wurde explizit in die Rolle eines
ablehnenden Gutachters versetzt mit dem Auftrag, das Paper abzulehnen. Sechs
Kritikpunkte wurden systematisch geprüft: Konfundierungen, Messprobleme,
Generalisierbarkeit. Punkte mit hoher Bewertung mussten vor dem Weiterschreiben
behoben werden, entweder durch zusätzliche Analysen oder Diskussionsteil-Ergänzungen.
Eine feindliche Rolle wurde gewählt, weil LLMs bei der Selbstprüfung eigene Fehler
übersehen \cite{badscientist2025}.

Die statistischen Anforderungen wurden in v3 erheblich verschärft. Zum Einsatz kamen
Boot\-strap-Kon\-fi\-denz\-inter\-valle, eine Korrektur für multiples Testen (BH-FDR),
eine vorab registrierte Mediationsanalyse sowie eine Multiversumsanalyse mit fünf
unterschiedlichen Spezifikationen. Für die Literaturrecherche folgte ein dokumentiertes PRISMA-lite-Protokoll,
ein vereinfachtes Screening-Verfahren, das Ein- und Ausschlussentscheidungen
transparent nachvollziehbar macht (36~Kandidaten gesichtet, 14~eingeschlossen).

\clearpage

\section{Inhaltliche Befunde}

Jede Pipeline-Version generierte autonom eine eigene Forschungsfrage und führte
die entsprechende statistische Modellierung durch.

Die \textbf{Baseline} untersuchte, ob Coding-Erfahrung mit Vertrauen in
KI-Ausgaben zusammenhängt. Es zeigte sich ein statistisch nachweisbarer, jedoch praktisch vernachlässigbarer
Effekt ($r_s = 0{,}10$). Der Agent kommunizierte diese geringe Effektstärke korrekt
und verzichtete auf eine Überinterpretation.

\textbf{Pipeline v1} fragte, ob wahrgenommene KI-Jobbedrohung die
Jobzufriedenheit stärker vorhersagt als das Jahresgehalt. Nach Listwise Deletion
($N = 10{,}112$) zeigte eine OLS-Regression (Likert-Skalen als quasi-metrisch
behandelt, in der Sozialforschung üblich), dass AIThreat der stärkste
standardisierte Prädiktor war, sogar stärker als das log-transformierte
Jahresgehalt. Bivariat lag der Effekt bei $d = -0{,}36$. Zudem zeigte sich, dass 81{,}6\,\% derjenigen, die KI als Bedrohung wahrnahmen,
KI-Tools dennoch aktiv nutzten. Wahrgenommene Bedrohung und tatsächliche Nutzung
schließen sich damit nicht aus.

\textbf{Pipeline v2} verglich AIThreat direkt mit dem tatsächlichen
KI-Nutzungsstatus (AISelect) als Prädiktoren ($N = 17{,}696$; größer als v1,
da das Jahresgehalt mit hohem Missing-Anteil nicht mehr als Kontrollvariable
einbezogen wurde). AIThreat zeigte einen substanziellen Effekt auf Jobzufriedenheit
($d = -0{,}33$), während AISelect praktisch keine Rolle spielte ($d = 0{,}03$,
nicht signifikant). Einstellung zu KI und tatsächliche Nutzung weisen damit eine divergente Wirkung
auf die Jobzufriedenheit auf.

\textbf{Pipeline v3} prüfte die in Stage~0b formulierte Hypothese: Moderiert
Vertrauen in KI-Genauigkeit (AIAcc) den negativen Effekt von AIThreat auf
Jobzufriedenheit ($N = 17{,}670$, nur aktive KI-Nutzer)? Das
Stress-Appraisal-Modell nach Lazarus und Folkman \cite{lazarus1984} diente als
theoretischer Hintergrund. Die Moderation ließ sich nicht nachweisen. Über alle
fünf Multiversum-Spezifikationen war der Interaktionsterm nicht signifikant
($f^2 < 0{,}001$). Weil die Hypothese vorab in Stage~0b festgelegt wurde, ist
dieses Null-Ergebnis eine valide wissenschaftliche Aussage.
Der AIThreat-Haupteffekt replizierte stabil ($d = -0{,}303$,
Bootstrap-KI $[-0{,}350;\, -0{,}250]$).

Über alle drei Pipelines hinweg zeigt sich der AIThreat-Effekt auf Jobzufriedenheit
konsistent: $d$ liegt zwischen $-0{,}30$ und $-0{,}36$, zuletzt unter
Pre-Registrierungsbedingungen. Einschränkend ist anzumerken, dass v2 und v3 auf dieselbe Ausgangsvariable wie
v1 zurückgriffen. Die Konsistenz ist damit teilweise durch das Design bedingt und stellt keine
vollständig unabhängige Replikation dar.

\section{Qualit\"atsevaluation}

Die Bewertungsmatrix wurde speziell für diesen Vergleich entwickelt. Sie bewertet
nur Kriterien, die von Pipelineentscheidungen abhängen. Sprache oder Struktur
blieben außen vor, da sie bei Ausgaben desselben Modells kaum variieren. Die
Scoring-Entscheidungen folgen ausschließlich den im Anhang definierten, regelbasierten
Ankerpunkten; es handelt sich um eine prozedurale, keine freie Beurteilung.
\clearpage
\begin{table}[h]
\centering
\caption{Bewertungsmatrix: pipeline-orientierte Kriterien (je 5 Punkte)}
\label{tab:matrix}
\small
\begin{tabular}{lcccc}
\toprule
\textbf{Kriterium} & \textbf{Baseline} & \textbf{v1} & \textbf{v2} & \textbf{v3} \\
\midrule
Literaturbasis \& Zitationsqualität & 2 & 3 & 4 & \textbf{5} \\
Methodische Robustheit & 2 & 3 & 4 & \textbf{5} \\
Wissenschaftliche Integrität & 1 & 2 & 3 & \textbf{5} \\
Inhaltliche Erkenntnistiefe & 3 & 4 & 3 & 4 \\
Verifikation \& Reproduzierbarkeit & 1 & 4 & 4 & \textbf{5} \\
\midrule
\textbf{Gesamt} & \textbf{9/25} & \textbf{16/25} & \textbf{18/25} & \textbf{24/25} \\
\bottomrule
\end{tabular}
\end{table}

Die vollständige Scoring-Rubrik mit Ankerpunkten findet sich in
Anhang~\ref{app:rubrik}.

v2 liegt mit 18/25 zwei Punkte vor v1. Die Differenz resultiert aus zwei Kriterien: Bei
\emph{Literaturbasis} schneidet v2 besser ab, weil der zusätzliche Critic-Check
tiefere Verifikation bietet als der reine Fetch-MCP-Abruf in v1. Bei
\emph{Wissenschaftliche Integrität} rechtfertigen die dokumentierten
SQL-Nachrechnungen und PAP-Abweichungsprotokolle einen höheren Score als die
formalen QC-Vorgaben in v1, die ohne aktive Fehlerkorrektur blieben.

Der größte Sprung liegt zwischen v2 und v3. \emph{Wissenschaftliche Integrität}
steigt von 3 auf 5 Punkte.
Pre-Registrierung und Adversarial Critic schaffen strukturelle Absicherungen,
die v2 noch nicht hatte.

\emph{Inhaltliche Erkenntnistiefe} ist in v3 nicht höher als in v1 (je 4 Punkte),
obwohl v3 methodisch weiterentwickelt ist. v3 produziert ein Null-Ergebnis, das
durch Pre-Registrierung als valide wissenschaftliche Aussage gesichert ist. Der
Score bewertet den Erkenntnisgewinn aus den Daten, nicht die methodische Qualität
des Ergebnisses selbst.

\section{Fazit und Ausblick}

\subsection{Fazit}

Die Ergebnisse bestätigen die Alternativhypothese: Eine strukturierte Pipeline
verbessert die wissenschaftliche Qualität autonomer Paper-Generierung messbar.
Auffällig ist dabei die ungleichmäßige Verteilung dieser Verbesserung. Eine
bloße Erhöhung der Werkzeuganzahl erbringt keinen nennenswerten Mehrwert. Den
Unterschied machen gezielte Eingriffe gegen spezifische Schwachstellen: die
STDOUT-Pflicht gegen Fake-Compliance, das Effect-Size-Gate gegen triviale Effekte,
die Pre-Registrierung gegen HARKing und der Adversarial Critic gegen unkritische
Selbstprüfung.

Dabei lässt sich ein konsistentes Muster erkennen. Textuelle Vorgaben erweisen sich
als deutlich weniger zuverlässig als Mechanismen, die direkt im Ausführungscode
verankert sind. Die Likert-Rebellion verdeutlicht dies für Methodenvorgaben; auch
Qualitätspflichten bleiben wirkungslos, wenn sie lediglich im Prompt formuliert sind
und nicht strukturell durchgesetzt werden. Code-Assertions, JSON-Locks und
Blocking-Bedingungen erwiesen sich dort als wirksam, wo deklarative
Prompt-Anweisungen scheiterten.

Die Pipeline weist dennoch klare Grenzen auf. Als universelles Werkzeug für beliebige
Datensätze ist sie nicht konzipiert. Ihr Wert liegt in der systematischen Verifikation
und Strukturierung des Forschungsprozesses. Die Beurteilung wissenschaftlicher Relevanz
und die Interpretation von Ergebnissen verbleiben beim Menschen.

\subsection{Limitationen}

Die Methodik weist spezifische Einschränkungen auf. Die Beschränkung auf einen
einzigen Datensatz und ein Sprachmodell verhindert eine direkte Generalisierung
der identifizierten Designprinzipien auf andere Kontexte. Jede Pipeline-Version
wurde zudem nur einmal ausgeführt, obwohl LLMs bei identischen Prompts
unterschiedliche Ergebnisse produzieren können. Einige der beobachteten
Qualitätsunterschiede könnten daher stochastisch bedingt sein. Erst mehrere
Durchläufe pro Bedingung könnten dies ausschließen.

Darüber hinaus unterscheiden sich die verglichenen Bedingungen nicht nur in der
Pipeline-Architektur, sondern auch in der jeweiligen Forschungsfrage. Komplexität
der Pipeline und Komplexität der Fragestellung lassen sich damit nicht trennscharf
unterscheiden,
was die Interpretation der gemessenen Qualitätssteigerungen erschwert. Eine
unabhängige Bewertung durch Dritte fehlt bislang völlig. Da der Survey querschnittlich angelegt ist, lässt sich der AIThreat-Effekt nicht
kausal interpretieren.

\subsection{Ausblick}

Für die Weiterentwicklung ergeben sich drei Anknüpfungspunkte. Erstens: Validierung
der Pipeline auf anderen Datensätzen, um zu überprüfen, ob die Designprinzipien auch
dort greifen. Zweitens: Vergleich verschiedener Sprachmodelle, um zu klären, ob
Phänomene wie die Likert-Rebellion oder Fake-Compliance modellspezifisch sind oder
universelle Eigenschaften leistungsfähiger LLMs darstellen. Drittens: Einbezug
externer Gutachter zur Validierung der Bewertungsmatrix an unabhängigen Urteilen.

\bibliographystyle{IEEEtran}
\bibliography{references}

\clearpage
\appendix

\section{Scoring-Rubrik der Bewertungsmatrix}
\label{app:rubrik}

\begin{table}[h]
\centering
\caption{Scoring-Rubrik: Ankerpunkte je Kriterium}
\label{tab:rubrik}
\footnotesize
\renewcommand{\arraystretch}{0.85}
\begin{tabular}{p{2.8cm}p{1.3cm}p{9cm}}
\toprule
\textbf{Kriterium} & \textbf{Punkte} & \textbf{Beschreibung} \\
\midrule
Literaturbasis \& Zitationsqualität
  & 5 & PRISMA-Protokoll, alle Quellen per Fetch-MCP sowie Critic-Check mit Abstract-Match verifiziert \\
  & 4 & Fetch-MCP sowie Critic-Check, Abstract-Verifikation dokumentiert \\
  & 3 & Fetch-MCP mit Live-Abruf, Relevanz dokumentiert, kein Critic-Check \\
  & 2 & Teilweise verifiziert, einige Quellen aus Trainingswissen \\
  & 1 & Keine Live-Verifikation, Quellen aus Trainingswissen \\
\midrule
Methodische Robustheit
  & 5 & Bootstrap-CIs, Multiverse ($\geq$5 Spez.), FDR-Korrektur, Mediation, Effektgröße \\
  & 4 & Robustness-Checks, VIF, Effektgröße, Missing-Data-Transparenz \\
  & 3 & Korrekte Tests, Effektgröße, keine Robustness-Checks \\
  & 2 & Korrekte Tests, aber Lücken (kein Post-hoc, Näherung statt exaktem p, falsche ES-Formel) \\
  & 1 & Methodenfehler oder nicht nachvollziehbar \\
\midrule
Wissenschaftliche Integrität
  & 5 & Pre-Registrierung, Adversarial Critic sowie explizite Null-Ergebnis-Dokumentation \\
  & 4 & Pre-Registrierung oder Adversarial Critic, aber nicht beides \\
  & 3 & Critic-Loops mit dokumentierten Korrekturen, PAP-Abweichung protokolliert \\
  & 2 & Formale QC-Vorgaben ohne aktive Fehlerkorrektur \\
  & 1 & Keine Qualitätskontrolle \\
\midrule
Inhaltliche Erkenntnistiefe
  & 5 & Neuer Befund, theoretische Einbettung, Implikationen, Multiverse-Bestätigung \\
  & 4 & Klarer Befund mit Kontextualisierung oder valides Null-Ergebnis mit epistemischer Absicherung \\
  & 3 & Befund kommuniziert, begrenzte Kontextualisierung \\
  & 2 & Befund vorhanden, kaum Interpretation \\
  & 1 & Kein klarer Befund \\
\midrule
Verifikation \& Reproduzierbarkeit
  & 5 & STDOUT-Belege für alle QC-Checks, SQL-Nachrechnung aller Kernzahlen \\
  & 4 & STDOUT-Belege für QC, SQL-Abfragen dokumentiert \\
  & 3 & Teilweise STDOUT-Belege, keine SQL-Nachrechnung \\
  & 2 & Code vorhanden, aber kein STDOUT-Nachweis \\
  & 1 & Keine Verifikation, kein STDOUT-Nachweis \\
\bottomrule
\end{tabular}
\end{table}

\end{document}
